% ============================================================================
% 1. INTRODUCTION
% ============================================================================
\section{Introduction}

Text-to-SQL has made rapid progress: state-of-the-art LLMs now approach or
match human-level accuracy on Spider~\cite{yu2018spider} and
BIRD~\cite{li2023bird}, which measure a model's capacity to translate a
natural-language question into a SQL query.  However, even the latest
benchmarks like Spider~2.0~\cite{lei2024spider2} show that enterprise-scale
SQLworkflows remain largely unsolved (best models~$\sim$17--21\%).

Data warehouses, on the other hand, present a different kind of problem.
A data engineer maintaining an enterprise warehouse does not primarily write
SQL.  Instead, they trace graph structure: lineage edges, foreign key (FK) paths,
disconnected schema silos.  Common questions include:

\begin{itemize}
  \item \emph{``If the \texttt{Person} table changes, which downstream tables
    are affected through data lineage?''} (impact analysis)
  \item \emph{``Which tables form disconnected silos in our schema graph?''}
    (connectivity analysis)
  \item \emph{``What is the shortest FK path between two tables?''}
    (routing)
\end{itemize}

Answering these questions demands breadth-first search (BFS), connected-component detection, and
multi-hop path enumeration.  No existing benchmark assesses LLMs on these capabilities
over real-world data warehouse schemas with heterogeneous
edge semantics.

Our evaluation reveals three principal findings:
(i)~all models achieve near-perfect scores on single-hop
structural queries but drop 30--40~percentage points on
compositional multi-hop tasks, exposing a systematic
\emph{structural reasoning ceiling};
(ii)~Tool-Use prompting closes the gap on topology-enumeration
subtypes but fails on compositional lineage-impact queries
that require chaining multiple graph algorithms; and
(iii)~obfuscating table names degrades performance by up to
32~pp on lineage tasks, confirming that models partially
rely on lexical cues rather than genuine graph traversal.



Our contributions are summarized as follows:

\begin{enumerate}
  \item \textbf{DW-Bench}: 1,046 schema-level questions
    across 5 data warehouse schemas (4 real-world, 1 synthetic), covering
    3 categories (lineage impact, schema routing, silo detection),
    13 subtypes, and 3 difficulty levels.
  \item \textbf{Six baselines} (Flat Text, Vector-RAG, Graph-Augmented,
    Tool-Use, ReAct-Code, Oracle), tested with Gemini~2.5 Flash,
    DeepSeek-V3, and Qwen2.5-72B.
  \item \textbf{An obfuscation protocol} that randomizes table names,
    enabling controlled evaluation of memorization versus genuine topology
    understanding.
  \item \textbf{Syn-Logistics}: A synthetic dataset with $n \geq 20$ per
    subtype.  It exposes a model-dependent structural gap: Gemini scores
    93--100\% on \texttt{hop\_count}/\texttt{count}, while DeepSeek drops to
    27--48\% on the same subtypes.
  \item \textbf{Open-source release} of all datasets, evaluation code, and
    baseline implementations.
\end{enumerate}
